\chapter{\altGerEng{Fazit}{Conclusion}}

Developing an energy-efficient Li-Fi receiver capable of harvesting power from the signal-light must address the wireless device’s limited battery life in \ac{VLC}. In this work, \ac{SLIPT} techniques are explored, implementing a prototype that uses the incoming LED light for both data reception and \ac{EH}. This thesis tried to deliver a battery-less proof-of-concept \ac{Li-Fi} receiver with \ac{SLIPT}. This thesis work represents a step towards self-sustaining optical wireless communication, reducing dependence on battery power for IoT and smart sensors.

The receiver’s communication module performed reliably under various test conditions. It transmitted data at 9600 to 19200 baud rates across distances of 5–15cm, with the STM32 instantly demodulating the information. Using a photodiode with amplification and filtering produced signal waveforms with much noise, but error rates were low enough to recover the transmitted text/data. Even higher baud rates(38400) were decodable at 15cm. This indicates that the design can support typical low-to-medium data rate \ac{Li-Fi} applications with short range. The \ac{EVM} managed to gather energy from the light but in limited amounts. With a dedicated mini solar panel as the collector, the receiver could convert the optical signal into electrical energy and gradually charge a 0.1F supercapacitor. In a low-power STOP mode, the power consumed by the circuit is reduced to 0.010W. The harvested energy extended the device’s operation, slowing the supercapacitor discharge. This confirms the fundamental premise that \ac{SLIPT} can work in a real \ac{Li-Fi} system, meaning the receiver processes data and gains energy from the medium.

This work also highlights the limitations and underperformance areas of the current design. A significant challenge was the extremely low current output from photodiodes in photovoltaic mode when illuminated by typical LED signals. The initial approach of using an array of nine photodiodes (for increased voltage) produced a voltage of 1.8V and a current of below 1\(\mu\)A, which was far too small to drive a boost converter. The attempt could not provide the necessary power for the receiver circuit. This is why another configuration combining a parallel and series of photodiodes was attempted. Even then, the energy harvested was modest. The small solar cell provided an improved current (10–15\(\mu\)A range under one LED at 5cm) but only at around 1V output. It initially required partially charging the supercapacitor to help bootstrap the harvesting circuit. In practice, it cannot start from a fully depleted state on harvested light alone – an external charge or illumination is needed. 

Another limitation observed was the trade-off between illumination intensity and communication fidelity. At very short distances, the photodiode and amplifier occasionally became saturated or behaved non-linearly, which made high-speed data transmission less reliable at close range than at a slightly greater distance. Nevertheless, using multiple LEDs to increase the harvested power also introduced interference and differences in timing that affected the received signal, as evident by the eye diagram. 

The experiments with a strong 60W incandescent lamp showed the gap between typical and ideal conditions. Under a concentrated 60W light source, the solar panel could fully charge the supercapacitor to 5.1V in under 3 minutes and even maintain a regulated output for a short while. Under the same 60W light, the photodiode array improved its performance (charging above 3.2V after 40 minutes) but could not sustain a load for long. The current design’s underperformance is mainly due to insufficient optical power and conversion efficiency and inability to design the \ac{EVM} circuit properly.

These findings indicate that the current receiver design operates at extremely tight power and signal margins. The \ac{EH} efficiency is low, and the system’s ability to self-sustain is constrained by the fundamental limits of photodiodes and the power needs of even low-power electronics. Photodiodes, especially when unbiased, do not efficiently convert light to electrical energy, meaning their responsivity and surface area limit the current. Using a compact photodiode array (for its fast response) meant sacrificing significant photocurrent compared to a larger-area solar cell. Conversely, while providing more current, the solar panel had a much lower bandwidth and could not be used to receive high-speed data. This fundamental trade-off concludes that devices optimized for power (large area, high capacitance) are not optimal for data and vice versa. Moreover, the power management circuit has an inherent threshold – it needs a certain minimum input power to start up. This was mitigated by pre-charging the storage or using a small hysteresis that allows intermittent output when the capacitor reaches 3.2V. Still, it could not maintain the voltage supply to the load.  


In conclusion, this thesis demonstrates that it is possible to simultaneously receive data and harvest power from the same light source. The insights gained from the issues faced include the need for better photodetectors, more efficient power management, and possibly revised \ac{SLIPT} architectures. By implementing the suggested improvements in \ac{EH} efficiency and optimizing the receiver circuit design, it is possible to build real batteryless \ac{Li-Fi} systems.

